%! Setting Up a Test for the Difference of Two Population Means — Unit 7, Lesson 7.8 — Warm-Up
\documentclass[10pt]{article}
\usepackage{apstats-article}
\usepackage{apstats-boxes}
\newcommand{\TallMath}[1]{$\displaystyle #1\rule[-1.4em]{0pt}{3.2em}$}

\begin{document}

\pageheader{Unit 7, Lesson 7.8}{Warm-Up}
\namedateperiod

\vspace{0.10in}

\begin{notesbox}{Spiral Review --- One-Sample $t$-Tests \& Two-Sample Conditions (Lessons 7.4, 7.6)}

\textbf{1.} A nutritionist believes the mean daily fiber intake of adults in a city is
\emph{less} than the recommended 25 g. She randomly samples 30 adults; $\sigma$ is unknown.

\begin{enumerate}[label=\alph*.]
  \item Name the appropriate inference method. \blank{5cm} \quad $df = $ \blank{1.5cm}
  \item State $H_0$ and $H_a$ using the population parameter.

        $H_0$: \blank{4cm} \qquad $H_a$: \blank{4cm}
  \item Is the Normal condition met? \blank{1.5cm} \; Explain briefly: \writeline
\end{enumerate}

\vspace{0.08in}
\textbf{2.} A researcher plans to compare mean exam scores for two independently selected
groups of students: those taught with Method A ($n_1 = 24$) and Method B ($n_2 = 28$).
Both populations are unknown.

\begin{enumerate}[label=\alph*.]
  \item Which conditions must be verified before running a two-sample $t$-test? Circle all that apply.

        \quad (A) Random for each group \quad
        (B) $n_1 + n_2 \ge 30$ \quad
        (C) 10\% rule for each group \quad
        (D) Both distributions checked when $n < 30$

  \item Define $\mu_1$ and $\mu_2$ in context.

        $\mu_1 = $ \blank{5.5cm} \qquad $\mu_2 = $ \blank{5.5cm}
\end{enumerate}

\end{notesbox}

\end{document}
